\documentclass[11pt]{article}
\usepackage[margin=1in]{geometry}
\usepackage{amsmath,amssymb,amsthm}
\usepackage[T1]{fontenc}
\usepackage{lmodern}
\usepackage{xcolor}
\usepackage{listings}
\usepackage{booktabs}
\usepackage[hidelinks]{hyperref}
\usepackage{url}

\newtheorem{theorem}{Theorem}
\newtheorem{lemma}[theorem]{Lemma}
\newtheorem{corollary}[theorem]{Corollary}
\newtheorem{proposition}[theorem]{Proposition}
\theoremstyle{remark}
\newtheorem{remark}[theorem]{Remark}

\newcommand{\lean}[1]{\texttt{#1}}
\newcommand{\br}[1]{[#1]}
\newcommand{\PP}{\mathbb{P}}
\newcommand{\repo}{https://github.com/jpbald93/pascal-hexagon}
% version-2 commits: pascal-hexagon=be35ffaa75a324131a710d3287b3d703408ad4a0
\newcommand{\commit}{726da4d}
\newcommand{\src}[2]{\href{\repo/blob/\commit/PascalHexagon/#1\##2}{\texttt{#1}}}

\lstdefinelanguage{lean}{
  morekeywords={theorem,lemma,def,example,by,fun,have,show,exact,intro,namespace,end,import,private,variable},
  sensitive=true, morecomment=[l]{--}, morecomment=[s]{/-}{-/}, morestring=[b]"}
\lstset{language=lean, basicstyle=\ttfamily\small, keywordstyle=\bfseries,
  columns=fullflexible, keepspaces=true, breaklines=true, frame=single, framerule=0.3pt,
  literate={ℙ}{{$\mathbb{P}$}}1 {≠}{{$\neq$}}1 {↔}{{$\leftrightarrow$}}2 {→}{{$\to$}}1
           {⨯₃}{{$\times_3$}}2 {⬝ᵥ}{{$\cdot_v$}}2 {₁}{{$_1$}}1 {₂}{{$_2$}}1 {₃}{{$_3$}}1
           {∀}{{$\forall$}}1 {∧}{{$\wedge$}}1 {∨}{{$\vee$}}1 {'}{{$'$}}1}

\title{Pascal's hexagon theorem over an arbitrary field,\\ formalised in Lean~4}
\author{Josh Bald\\ \small Independent researcher, Ontario, Canada\\ \small\texttt{jpbald93@gmail.com}}
\date{28 September 2026; version 2, 22 October 2026 (adds Section~\ref{sec:l1415})}

\emergencystretch=3em
\begin{document}
\maketitle

\begin{abstract}
We formalise Pascal's hexagon theorem in Lean~4 with Mathlib. The statement covers an arbitrary
conic over an arbitrary field: six points on the zero set of a nonzero ternary quadratic form. It
places no restriction on the characteristic and allows degenerate conics. The only hypotheses are
the ones needed for the three intersection points of opposite sides to be defined. Pappus's
theorem follows as the case of a line pair, and Brianchon's dual-conic formulation as the projective dual. The formal statement uses Mathlib's projective plane
$\PP(K^3)$, its cross-product construction of lines and intersection points, and its notion of
collinearity. We also prove the lemmas a reader needs to check that this statement means what it
says. The proof rests on a single polynomial identity, valid over any commutative ring: the triple
product of the three Pascal points is, up to sign, the $6\times6$ determinant of the degree-two
Veronese images of the six vertices. The mathematics is classical. As far as we know, this is the first
Lean formalisation, without \lean{sorry} or added axioms, for an arbitrary nonzero ternary
quadratic form over an arbitrary field. It extends the real projective statements formalised in
HOL~Light and Mizar to arbitrary fields. Those statements assume 17 or 18 non-collinearity
conditions on the vertices; we replace them by distinctness of consecutive vertices and of
opposite sides. The development has 32 named theorems in about 670 lines. A
checking script, tested against deliberately tampered copies of the source, confirms that each of
them uses only Lean's standard axioms. Version~2 adds Mathlib-only restatements with explicit
non-degenerate witnesses, and a comparison with independently written formalisations of the textbook
statements. That comparison proves, outside the main development, Brianchon's theorem for tangent lines
in characteristic $\ne2$, the five-point (tangent) form of Pascal's theorem and its converse.
\end{abstract}

\section{Introduction}

Pascal's theorem (1640) is the statement that the three points where opposite sides of a hexagon
inscribed in a conic meet are collinear. It is number~28 on Wiedijk's list of 100
theorems~\cite{Wiedijk}. The list records formalisations in HOL~Light, Mizar and Rocq. Mathlib's
copy of the list has no Lean declaration for it~\cite{Mathlib100}.

We prove the following.

\begin{theorem}[Pascal]\label{thm:pascal}
Let $K$ be a field and $Q$ a nonzero quadratic form on $K^3$. Let $A,B,C,D,E,F$ be points of the
projective plane $\PP^2(K)$ with $Q(A)=\dots=Q(F)=0$. Suppose consecutive vertices are distinct,
$A\neq B$, $B\neq C$, $\dots$, $F\neq A$, and opposite sides are distinct lines,
$AB\neq DE$, $BC\neq EF$, $CD\neq FA$. Then the points $AB\cap DE$, $BC\cap EF$ and $CD\cap FA$ are
collinear.
\end{theorem}

No assumption is made on the characteristic of $K$, on the nondegeneracy of $Q$, or on
non-consecutive vertices. The hypotheses are exactly what is needed for the six sides to be lines
and for the three intersection points to exist. If $Q=\ell_1\ell_2$ is a product of two linear forms,
Theorem~\ref{thm:pascal} is Pappus's theorem (Corollary~\ref{cor:pappus}).

\paragraph{Contributions.}
\begin{itemize}
\item A Lean~4 formalisation of Theorem~\ref{thm:pascal} (\lean{PascalHexagon.pascal}). It is
stated with Mathlib's \lean{Projectivization}, \lean{Projectivization.cross} and
\lean{Projectivization.IsCollinear}.
\item A vector form over any integral domain with no distinctness hypotheses at all
(Theorem~\ref{thm:vector}), and the underlying identity over any commutative ring
(Proposition~\ref{prop:identity}).
\item Lemmas that justify the formal statement (Section~\ref{sec:correct}). They show that
\lean{cross A B} is the unique line through $A$ and $B$, that \lean{cross L M} is the unique point
on $L$ and $M$, that collinearity is equivalent to the vanishing of a determinant, and that lying on
the conic does not depend on the chosen representative. We also give a version whose conclusion does
not mention \lean{cross} at all.
\item Pappus's theorem as a corollary, and Brianchon's dual-conic formulation as the projective
dual (Section~\ref{sec:brianchon}).
\item A tamper-tested checking script that confirms the axiom footprint of every named theorem
(Section~\ref{sec:gate}).
\item In version~2 (Section~\ref{sec:l1415}): Mathlib-only restatements of the main results, with
explicit non-degenerate witnesses for statements with hypotheses, and a comparison with a blind formalisation of the textbook
statements. The comparison proves all 20 textbook statements, among them Brianchon's theorem for
tangent lines (characteristic $\ne2$), the tangent form of Pascal's theorem and the
Braikenridge--Maclaurin converse.
\end{itemize}

\paragraph{Availability.}
The code is at \url{\repo}. Line references in this paper are to commit \texttt{\commit}. The
toolchain is Lean \texttt{v4.33.1}, and Mathlib is pinned in \texttt{lake-manifest.json}. The code is
released under Apache~2.0. Version~2 of this paper adds the files \lean{PascalHexagon/L14.lean} and
\lean{PascalHexagon/L15Classical.lean} (Section~\ref{sec:l1415}) in a new commit, \texttt{be35ffaa75a324131a710d3287b3d703408ad4a0}.
In that commit the root import file and \texttt{gate.sh} gain entries for them, and \texttt{gate.sh}
also gets a portability fix (Section~\ref{sec:l1415}). The earlier statements and files are
unchanged, so the line references hold there as well.
Supporting material (the Layer-14 and Layer-15 reports, the blind writer's statements and report, the transcript audit, and the reviews of this version with the fixes made in response) is in the folder \texttt{Pascal-Brianchon} of the repository \url{https://github.com/jpbald93/LEAN-verification}.

\begin{table}[t]
\centering\small
\begin{tabular}{@{}lll@{}}
\toprule
Lean name (namespace \lean{PascalHexagon}) & In this paper & Source \\ \midrule
\lean{pascalDet\_eq\_pascalBracket} & Proposition~\ref{prop:identity}(a) & \src{Bracket.lean}{L11} \\
\lean{det\_verMatrix\_eq\_neg\_pascalBracket} & Proposition~\ref{prop:identity}(b) & \src{Veronese.lean}{L12} \\
\lean{pascalDet\_eq\_zero\_of\_quadEval} & Lemma~\ref{lem:poly} & \src{Conic.lean}{L28} \\
\lean{apply\_eq\_quadEval} & Lemma~\ref{lem:coeffs} & \src{Conic.lean}{L49} \\
\lean{coeffs\_ne\_zero} & Lemma~\ref{lem:coeffs} & \src{Conic.lean}{L75} \\
\lean{pascal\_quadraticMap} & Theorem~\ref{thm:vector} & \src{Conic.lean}{L86} \\
\lean{pascal\_matrix} & Remark~\ref{rem:matrix} & \src{Conic.lean}{L103} \\
\lean{pascalDet\_standardConic} & Remark~\ref{rem:standard} & \src{StandardConic.lean}{L17} \\
\lean{isCollinear\_mk\_iff} & Lemma~\ref{lem:collinear} & \src{Correctness.lean}{L64} \\
\lean{mem\_line\_cross} & Lemma~\ref{lem:incidence}(a) & \src{Correctness.lean}{L91} \\
\lean{cross\_mem\_lines} & Lemma~\ref{lem:incidence}(a) & \src{Correctness.lean}{L96} \\
\lean{eq\_cross\_of\_orthogonal} & Lemma~\ref{lem:incidence}(b) & \src{Correctness.lean}{L102} \\
\lean{eq\_cross\_of\_orthogonal'} & Lemma~\ref{lem:incidence}(b), dual & \src{Correctness.lean}{L112} \\
\lean{cross\_cross\_spec} & Lemma~\ref{lem:incidence}(c) & \src{Correctness.lean}{L128} \\
\lean{onConic\_mk\_iff} & Lemma~\ref{lem:onconic} & \src{Correctness.lean}{L146} \\
\lean{pascal} & Theorem~\ref{thm:pascal} & \src{Projective.lean}{L92} \\
\lean{pascal'} & Theorem~\ref{thm:pascal}, with \lean{OnConic} & \src{Correctness.lean}{L155} \\
\lean{pascal\_of\_isCollinear} & Theorem~\ref{thm:pascal-col} & \src{Correctness.lean}{L168} \\
\lean{linePair\_ne\_zero} & Corollary~\ref{cor:pappus}, proof & \src{Correctness.lean}{L204} \\
\lean{pappus} & Corollary~\ref{cor:pappus} & \src{Correctness.lean}{L236} \\
\lean{brianchon} & Theorem~\ref{thm:brianchon} & \src{Brianchon.lean}{L38} \\
\lean{brianchon\_concurrent} & Theorem~\ref{thm:brianchon}, concurrency & \src{Brianchon.lean}{L61} \\
\bottomrule
\end{tabular}
\caption{The main formal statements. \texttt{gate.sh} checks these 22 theorems and 10 supporting
lemmas, 32 theorems in all.}
\label{tab:thms}
\end{table}

\section{Conventions and the Mathlib interface}\label{sec:conv}

We write vectors of $K^3$ as \lean{Fin 3} $\to$ \lean{K}. The bracket of three vectors is
$\br{u,v,w}=u\cdot(v\times w)=\det(u,v,w)$. Mathlib provides the cross product \lean{u} $\times_3$ \lean{v} and
the identity $\br{u,v,w}=\det(u,v,w)$ (\lean{triple\_product\_eq\_det}).

The projective plane is Mathlib's $\mathbb{P}$\,\lean{K (Fin 3} $\to$ \lean{K)}. Its points are nonzero vectors up to
nonzero scalars. \lean{P.rep} is a chosen representative, and \lean{mk K v hv} is the point of a
nonzero vector~$v$. Mathlib identifies lines with points of the same space via the dot product:
\lean{P.orthogonal L} means $P\cdot L=0$, which we read as ``$P$ lies on~$L$''. For distinct points
$A\ne B$, Mathlib's \lean{cross A B} is the point represented by $A\times B$. It serves both as the
line through two points and as the intersection point of two lines. Collinearity is Mathlib's
\lean{IsCollinear}: a set of points is collinear if it lies in a subspace of dimension at most~$2$.
Mathlib's definition of \lean{cross} requires \lean{[DecidableEq K]}, so the statements that mention
\lean{cross} carry this hypothesis. It is harmless in classical logic.

A conic is the zero set of a nonzero \lean{QuadraticForm K (Fin 3} $\to$ \lean{K)}. Here ``nonzero'' means
nonzero as a quadratic form, which is the same as having a nonzero coefficient vector
(Lemma~\ref{lem:coeffs}). We write \lean{OnConic Q P} for $Q(\lean{P.rep})=0$.

\section{The key identity}\label{sec:identity}

For vectors $A,\dots,F\in K^3$, define the \emph{Pascal determinant}
\[
\pi(A,\dots,F)=\br{(A\times B)\times(D\times E),\ (B\times C)\times(E\times F),\ (C\times D)\times(F\times A)}.
\]
If $A\ne B$ and $D\ne E$ as projective points and $AB\ne DE$, then $(A\times B)\times(D\times E)$
represents $AB\cap DE$, and similarly for the other two points. So Pascal's theorem says that
$\pi$ vanishes. Let $\nu(x,y,z)=(x^2,y^2,z^2,xy,yz,zx)$ be the degree-two Veronese map, and let
$V(A,\dots,F)$ be the $6\times6$ matrix with rows $\nu(A),\dots,\nu(F)$.

\begin{proposition}\label{prop:identity}
Over any commutative ring, for all $A,\dots,F\in K^3$:
\begin{enumerate}
\item[(a)] $\pi(A,\dots,F)=\beta(A,\dots,F)$, where $\beta$ is the bracket polynomial
\begin{multline*}
\beta=\br{ABE}\br{BCF}\br{CDA}\br{DEF}-\br{ABE}\br{BCF}\br{CDF}\br{DEA}\\
+\br{ABE}\br{BCE}\br{CDF}\br{DFA}-\br{ABD}\br{BCE}\br{CDF}\br{EFA};
\end{multline*}
\item[(b)] $\det V(A,\dots,F)=-\beta(A,\dots,F)$.
\end{enumerate}
Hence $\pi=-\det V$.
\end{proposition}

\begin{proof}
(a) Apply the vector triple product identity $(u\times v)\times w=(u\cdot w)v-(v\cdot w)u$ to each
of the three Pascal points. This gives, for instance,
\[
(A\times B)\times(D\times E)=\br{A,D,E}B-\br{B,D,E}A.
\]
Expanding the bracket of the three resulting vectors by multilinearity, and cancelling terms with a
repeated argument, gives an expression in brackets. Its equality with $\beta$ is a polynomial
identity in the $18$ coordinates. (b) This is a polynomial identity in the $18$ coordinates.
\end{proof}

Formally, (a) is proved by rewriting with Mathlib's \lean{cross\_cross\_eq\_smul\_sub\_smul} and
then calling \lean{ring} on the coordinate expansion. (b) is proved by Laplace expansion along the
first row, followed by \lean{ring}. Neither proof is pretty. Each takes about 80--95~seconds and
2.5~GB of memory to check. They are, however, short, and they need no hypotheses. Brackets for the
conic condition are classical (see e.g.~\cite{RichterGebert}). HOL~Light's formalisation also goes
through a bracket condition (Section~\ref{sec:related}).

\begin{lemma}\label{lem:poly}
Let $K$ be an integral domain and $c\in K^6$ nonzero. If $\nu(v)\cdot c=0$ for $v=A,\dots,F$, then
$\pi(A,\dots,F)=0$.
\end{lemma}
\begin{proof}
$V(A,\dots,F)\,c=0$ with $c\ne0$, so $\det V=0$, and $\pi=-\det V=0$ by
Proposition~\ref{prop:identity}.
\end{proof}

\begin{lemma}\label{lem:coeffs}
For every quadratic map $Q$ on $K^3$ over a commutative ring, $Q(v)=\nu(v)\cdot c_Q$, where
\[
c_Q=\bigl(Q(e_1),Q(e_2),Q(e_3),\,b_Q(e_1,e_2),\,b_Q(e_2,e_3),\,b_Q(e_3,e_1)\bigr)
\]
and $b_Q(x,y)=Q(x+y)-Q(x)-Q(y)$ is the polar form. If $Q\ne0$, then $c_Q\ne0$.
\end{lemma}

\begin{theorem}[vector form]\label{thm:vector}
Let $K$ be an integral domain and $Q\ne0$ a quadratic map on $K^3$. If $A,\dots,F\in K^3$ satisfy
$Q(A)=\dots=Q(F)=0$, then $\pi(A,\dots,F)=0$.
\end{theorem}

\begin{proof}
Lemmas~\ref{lem:coeffs} and~\ref{lem:poly}.
\end{proof}

\begin{lstlisting}
theorem pascal_quadraticMap [IsDomain K] {Q : QuadraticMap K (Fin 3 → K) K} (hQ : Q ≠ 0)
    {A B C D E F : Fin 3 → K}
    (hA : Q A = 0) (hB : Q B = 0) (hC : Q C = 0) (hD : Q D = 0) (hE : Q E = 0) (hF : Q F = 0) :
    ((A ⨯₃ B) ⨯₃ (D ⨯₃ E)) ⬝ᵥ ((B ⨯₃ C) ⨯₃ (E ⨯₃ F)) ⨯₃ ((C ⨯₃ D) ⨯₃ (F ⨯₃ A)) = 0
\end{lstlisting}

Theorem~\ref{thm:vector} has no distinctness hypotheses at all. When the configuration is
degenerate, for example $A=B$, some Pascal point is the zero vector and the statement holds
trivially. The content of Pascal's theorem lies in the nondegenerate case, which the projective
form below isolates.

\begin{remark}\label{rem:matrix}
Many texts describe a conic by a symmetric matrix $M$ with $v^{\mathsf T}Mv=0$. The formal version
\lean{pascal\_matrix} is stated over a field $K$ with $2\ne0$: if $M\ne0$ is symmetric and
$v^{\mathsf T}Mv=0$ for $v=A,\dots,F$, then $\pi(A,\dots,F)=0$. The hypothesis $2\ne0$ is needed. In characteristic~$2$, $v^{\mathsf T}Mv$
sees only the diagonal of~$M$, so a nonzero symmetric matrix can give the zero form. This is why
we state the main results with quadratic forms.
\end{remark}

\begin{remark}\label{rem:standard}
The development also contains a direct proof for the standard parametrised conic $xz=y^2$, with
points $(s^2,st,t^2)$ (\lean{pascalDet\_standardConic}). It holds for all parameters $s,t$ over any
commutative ring, and it is not used in the main proof.
\end{remark}

\section{The projective statement}\label{sec:proj}

\begin{lstlisting}
theorem pascal {Q : QuadraticForm K (Fin 3 → K)} (hQ : Q ≠ 0)
    {A B C D E F : ℙ K (Fin 3 → K)}
    (hQA : Q A.rep = 0) (hQB : Q B.rep = 0) (hQC : Q C.rep = 0) (hQD : Q D.rep = 0)
    (hQE : Q E.rep = 0) (hQF : Q F.rep = 0)
    (hAB : A ≠ B) (hBC : B ≠ C) (hCD : C ≠ D) (hDE : D ≠ E) (hEF : E ≠ F) (hFA : F ≠ A)
    (h₁ : cross A B ≠ cross D E) (h₂ : cross B C ≠ cross E F) (h₃ : cross C D ≠ cross F A) :
    IsCollinear ({cross (cross A B) (cross D E), cross (cross B C) (cross E F),
      cross (cross C D) (cross F A)} : Set (ℙ K (Fin 3 → K)))
\end{lstlisting}

The field $K$ carries \lean{[Field K] [DecidableEq K]}. The proof unfolds \lean{cross} on
representatives. Theorem~\ref{thm:vector} then shows that the triple product of the representatives
of the three points vanishes, and Lemma~\ref{lem:collinear}(b) turns this into collinearity.

\begin{lemma}\label{lem:collinear}
For nonzero $u,v,w\in K^3$:
\begin{enumerate}
\item[(a)] if $\{[u],[v],[w]\}$ is collinear, then $\br{u,v,w}=0$;
\item[(b)] if $\br{u,v,w}=0$, then $\{[u],[v],[w]\}$ is collinear.
\end{enumerate}
\end{lemma}

Part (b) is \lean{isCollinear\_mk\_of\_triple\_product\_eq\_zero}. It is used in the proof of
\lean{pascal}. Part (a) is its converse, and the two together give \lean{isCollinear\_mk\_iff}.
Part (a) is not needed for Pascal's theorem. We include it so that ``collinear'' in the formal
statement can be read as the familiar determinant condition.

\section{Does the statement mean what it says?}\label{sec:correct}

A reader of \lean{pascal} must trust two definitions: that \lean{cross A B} is the line $AB$, and
that \lean{cross (cross A B) (cross D E)} is the point $AB\cap DE$. Mathlib proves that these
points are incident (\lean{cross\_orthogonal\_left} and its variants). We restate this and add
uniqueness.

\begin{lemma}\label{lem:incidence}
Let $A,B,L,M,P\in\PP^2(K)$.
\begin{enumerate}
\item[(a)] If $A\neq B$, then $A$ and $B$ lie on \lean{cross A B}. If $L\neq M$, then
\lean{cross L M} lies on $L$ and on~$M$.
\item[(b)] If $L\neq M$ and $P$ lies on both $L$ and $M$, then $P=\lean{cross L M}$. Dually, if
$A\ne B$ and both lie on a line~$L$, then $L=\lean{cross A B}$.
\item[(c)] If $A\neq B$, $D\ne E$ and $\lean{cross A B}\ne\lean{cross D E}$, then the point
\[\lean{cross (cross A B) (cross D E)}\]
is collinear with $A,B$ and with $D,E$, and it is the only
such point.
\end{enumerate}
\end{lemma}

Part (c) connects Mathlib's notion of collinearity with the cross-product construction. It gives a
version of Pascal's theorem whose conclusion does not mention \lean{cross} at all.

\begin{theorem}\label{thm:pascal-col}
Under the hypotheses of Theorem~\ref{thm:pascal}, let $P,R,S$ be points such that $P$ is
collinear with $A,B$ and with $D,E$, the point $R$ with $B,C$ and with $E,F$, and $S$ with $C,D$ and with $F,A$.
Then $\{P,R,S\}$ is collinear.
\end{theorem}

This is the form closest to the informal statement, and to the statements in HOL~Light and Mizar
(Section~\ref{sec:related}). The only occurrences of \lean{cross} are the hypotheses that opposite
sides are distinct. These hypotheses cannot be dropped. If $AB=DE$, every point of that line is
collinear with $A,B$ and with $D,E$, so $P$ is no longer determined. For example, over
$\mathbb{Q}$ take the six affine points $A=(0,0)$, $B=(1,0)$, $C=(0,1)$, $D=(2,0)$, $E=(3,0)$,
$F=(0,2)$ on the line pair $xy=0$. Here $AB=DE$, but $BC\ne EF$ and $CD\ne FA$. Their intersection
points are $R=(-3,4)$ and $S=(0,1)$. With $P=A$, the points $P,R,S$ are not collinear.

\begin{lemma}\label{lem:onconic}
For nonzero $v\in K^3$, \lean{OnConic Q (mk K v hv)} holds if and only if $Q(v)=0$.
\end{lemma}
\begin{proof}
Every representative has the form $\lambda v$ with $\lambda\ne0$, and $Q(\lambda v)=\lambda^2Q(v)$.
\end{proof}

So the conic hypotheses \lean{Q A.rep = 0} of \lean{pascal} do not depend on the representative.
\lean{pascal'} restates the theorem with \lean{OnConic}.

\begin{corollary}[Pappus]\label{cor:pappus}
Let $a,b$ be lines, let $A,C,E$ lie on $a$ and $B,D,F$ lie on~$b$. If consecutive vertices are
distinct and opposite sides are distinct lines, then $AB\cap DE$, $BC\cap EF$ and $CD\cap FA$ are
collinear.
\end{corollary}
\begin{proof}
The quadratic form $Q(x)=(a\cdot x)(b\cdot x)$ vanishes on all six points. It is nonzero, since
over a field there is a vector $z$ with $a\cdot z\ne0$ and $b\cdot z\ne0$. Apply
Theorem~\ref{thm:pascal}.
\end{proof}

The corollary allows $a=b$. In that case all six points lie on one line, and the hypotheses force
every side to be that line, so the hypothesis $AB\ne DE$ fails. The corollary is vacuous there, as
expected.

\section{Brianchon's theorem}\label{sec:brianchon}

In $\PP^2(K)$ with Mathlib's conventions, lines and points are the same objects. A point $P$ lies
on a line $L$ when $P\cdot L=0$, and \lean{cross} computes both the line through two points and the
intersection of two lines. Projective duality is therefore literal. Over a field of characteristic
different from two, the lines tangent to a nonsingular conic form a nonsingular conic in the dual
plane, the \emph{dual conic}. Pascal's theorem, applied in the dual plane, is Brianchon's theorem.

\begin{theorem}[Brianchon]\label{thm:brianchon}
Let $Q$ be a nonzero quadratic form on $K^3$, and let $a,b,c,d,e,f$ be lines lying on the dual
conic $Q=0$. Suppose consecutive sides are distinct and opposite vertices are distinct points:
$a\cap b\ne d\cap e$, $b\cap c\ne e\cap f$, $c\cap d\ne f\cap a$. Then the three diagonals joining
opposite vertices are collinear as points of the dual plane. Geometrically this is equivalent to
concurrency. Our explicit common-point corollary additionally assumes that the first two diagonals
are distinct (the Layer-14 restatement drops this assumption, Section~\ref{sec:l1415}).
\end{theorem}

\begin{proof}
The first statement is \lean{pascal'} applied to the lines $a,\dots,f$. For the second, if three
collinear points $L\ne M, N$ of the dual plane are given, then \lean{cross L M} lies on all three
lines by Lemma~\ref{lem:incidence}.
\end{proof}

This extra assumption is needed only for the chosen cross-product witness.

The formal statement takes the dual conic as its hypothesis. It does not construct the dual conic
of a given conic, and it does not define tangency. For a nondegenerate conic over a field of
characteristic $\ne2$ with Gram matrix $M$, the dual conic has Gram matrix $M^{-1}$, and a line is
tangent to the conic exactly when it lies on the dual conic. We have not formalised this
correspondence, so Theorem~\ref{thm:brianchon} is Brianchon's theorem only under this standard
identification. Version~2 proves Brianchon's theorem for six distinct tangent lines of a
nondegenerate conic in characteristic $\ne2$ separately, through the polarity of the conic and
\lean{pascal'}, in the Layer-15 file (Section~\ref{sec:l1415}). That proof does not use
Theorem~\ref{thm:brianchon}, and the Gram-matrix description of the dual conic is still not
formalised. For arbitrary nonzero dual-plane quadratic forms, including degenerate forms and characteristic
two, the formal result is a dual-Pascal generalisation; no tangent interpretation is asserted.

\section{Related work}\label{sec:related}

\paragraph{HOL Light.} Harrison's \texttt{100/pascal.ml}~\cite{HOLLight} proves Pascal's theorem
for a projective conic over~$\mathbb{R}$ given by six not-all-zero coefficients. Its statement
names the three intersection points $x_7,x_8,x_9$ by collinearity conditions, as in our
Theorem~\ref{thm:pascal-col}. It also assumes that 18 of the 20 triples of vertices are \emph{not} collinear. Affine and circle versions follow; the circle version is the one listed
on Wiedijk's page. The proof also uses a bracket condition for six points on a conic, checked with
\texttt{REAL\_RING}. Our approach is therefore close to Harrison's. The differences lie in the
generality of the statement. We work over an arbitrary field instead of~$\mathbb{R}$. Our
hypotheses only require the three intersection points to exist. The 18 non-collinearity
conditions are stronger. For example, they require that no three consecutive vertices of the
hexagon are collinear. In HOL~Light's labelling the hexagon is $x_1x_5x_3x_4x_2x_6$, and the two
triples it does not exclude, $\{x_1,x_2,x_3\}$ and $\{x_4,x_5,x_6\}$, are the alternate
vertices. So HOL~Light's statement does cover the line-pair case of Pappus's theorem, since its
conics may be degenerate.

\paragraph{Mizar.} Coghetto's article \texttt{PASCAL}~\cite{Coghetto} proves Pascal's theorem in
the real projective plane, for a conic given by six real coefficients that are not all zero. The
hypotheses form a ``Pascal configuration'' with 17 non-collinearity conditions, similar to
HOL~Light's. With the same labelling as HOL~Light, it leaves unconstrained the two alternate-vertex triples and one further triple, $\{P_3,P_5,P_6\}$.

\paragraph{Rocq.} Magaud and Narboux~\cite{MagaudNarboux} work in synthetic projective geometry. Their
result \texttt{hexamy\_prop} assumes a strong form of Pappus's theorem and five inequalities
between lines. It shows that the ``hexamy'' property is preserved when the first two vertices are
swapped. A hexamy is a hexagon with six pairwise distinct vertices for which the Pascal conclusion
holds. It is a statement about
axiomatic projective planes, and does not refer to conics defined by equations.

\paragraph{Lean.} We know of three public Lean~4 developments, none in Mathlib.
\begin{itemize}
\item One proves the theorem for six distinct points on a circle in the Euclidean
plane~\cite{PlbyLean}, without \lean{sorry}.
\item One gives a projective real version whose main theorem depends on a declared
\lean{axiom}~\cite{LeanGenius}. Some special cases, such as the standard conic, are proved without it.
\item One is an incomplete development for irreducible conics over a field, which relies on a
configuration lemma proved with \lean{sorry}~\cite{AtlasLean}.
\end{itemize}
We searched GitHub and the web for others; the search is not exhaustive. Mathlib's list of 100
theorems has no entry for \#28 at the version we use~\cite{Mathlib100}.

\section{Verification}\label{sec:gate}

\texttt{gate.sh} builds the library and then writes its own temporary file. That file imports the
library and issues \lean{\#print axioms} for each of the 32 named theorems. The queries are
therefore generated by the gate and not read from the repository. The gate flattens Lean's wrapped
output lines. It requires exactly one report per named theorem, and accepts only the axioms
\lean{propext}, \lean{Classical.choice} and \lean{Quot.sound}. Before building, it rejects source
files that contain \lean{sorry}, \lean{admit}, \lean{native\_decide}, \lean{axiom}, \lean{\#eval},
\lean{run\_cmd}, \lean{set\_option}, \lean{macro}, \lean{elab}, \lean{syntax}, \lean{notation},
\lean{import Lean}, or related tokens. At commit \texttt{\commit} it prints
\texttt{PASS (32 theorems, standard axioms only)}. The script needs GNU tools. Its \texttt{mktemp
--suffix} option is GNU-only, so with the BSD tools of stock macOS it stops with
\texttt{FAIL: axiom report}.

\texttt{tests/tamper.sh} plants six kinds of defect in scratch copies of the source and checks that
the gate rejects each one:
\begin{itemize}
\item a \lean{sorry};
\item an \lean{axiom} split across two lines;
\item an \lean{\#eval} that prints a fake axiom report;
\item a custom macro;
\item an altered incidence hypothesis in \lean{pappus} that is incompatible with the unchanged
proof;
\item a genuine extra axiom with a long name. The report wraps over several lines, which tests
the parser. For this test only, the check for the \lean{axiom} token is switched off in the scratch copy.
\end{itemize}
All six are rejected.

\texttt{code/check\_pascal.py} checks the identities with exact rational arithmetic in 300 random
trials:
\begin{itemize}
\item random points on random conics, which give $\pi=0$;
\item the identity $\pi=-\det V$ on arbitrary points;
\item the parametrised standard conic;
\item line pairs.
\end{itemize}
As a control, the script prints $\pi$ for five perturbed configurations; all five values are nonzero.

The gate is a heuristic, not a sandbox. It trusts the pinned toolchain and Mathlib, and it checks
only the named theorems. The step that matters most, comparing the formal statements with
Theorem~\ref{thm:pascal}, cannot be automated. Section~\ref{sec:correct} is meant to make that
comparison easier. The author has checked each statement in Table~\ref{tab:thms} against its
informal counterpart by hand.

\subsection{Layers 14 and 15}\label{sec:l1415}

Version~2 adds two further statement checks. Both are files of the library in commit
\texttt{be35ffaa75a324131a710d3287b3d703408ad4a0}, \lean{PascalHexagon/L14.lean} and \lean{PascalHexagon/L15Classical.lean}
(namespace \lean{PascalHexagon.Classical}). The earlier files are unchanged; the root import file and
\texttt{gate.sh} gain the new entries. The Layer-15 comparison was first checked as a scratch file
and then moved into the library. Below, Layer-14 names are in namespace
\lean{PascalHexagon} and Layer-15 names in \lean{PascalHexagon.Classical}.

\paragraph{Layer 14: statements in Mathlib's language and non-vacuity.} The main statements of
Table~\ref{tab:thms}, except the proof aid \lean{linePair\_ne\_zero}, are restated with Mathlib-only types and proved from the library theorem through
bridges. The restatements unfold \lean{pascalDet}, \lean{quadEval}, \lean{ver}, \lean{coeffs} and
\lean{OnConic} into Mathlib's \lean{crossProduct}, dot product, \lean{Matrix.det},
\lean{QuadraticMap.polar} and \lean{Projectivization.rep}. Sixteen theorems restate 21 of the 22
names of Table~\ref{tab:thms}. One theorem covers \lean{pascal}, \lean{pascal'} and
\lean{pascal\_of\_isCollinear}, one covers \lean{brianchon} and \lean{brianchon\_concurrent}, and the
four names of Lemma~\ref{lem:incidence}(a,b) are grouped into two (existence and uniqueness of the
line through two points and of the point on two lines). The 22nd name,
\lean{linePair\_ne\_zero}, is not restated. It is a step in the proof of Corollary~\ref{cor:pappus},
stated in terms of the library's \lean{linePair}, and it is used in the proof of the restated Pappus
theorem. A mechanical check of the elaborated types found no \lean{PascalHexagon} constant. Three
restatements are slightly stronger in form:
\begin{itemize}
\item \lean{l14\_pascal} and \lean{l14\_pappus} do not mention \lean{cross} at all.
The hypothesis that opposite sides are distinct is written as
$\neg(\{A,B,D\}\text{ and }\{A,B,E\}\text{ both collinear})$, which for $A\ne B$, $D\ne E$ is
equivalent to $\lean{cross A B}\ne\lean{cross D E}$ (bridge \lean{l14b\_cross\_ne\_iff}).
\item \lean{l14\_brianchon} proves that the three diagonals have a common point
\emph{without} the hypothesis that the first two diagonals are distinct
(bridge \lean{l14b\_common\_point}). It still takes the dual conic as a hypothesis, as
in Section~\ref{sec:brianchon}.
\end{itemize}
Five witness theorems prove every hypothesis on explicit configurations over $\mathbb{Q}$ and that
the conclusion is non-degenerate:
\begin{center}\footnotesize
\begin{tabular}{@{}lp{0.62\textwidth}@{}}
\toprule
statement & witness and non-degeneracy proved\\
\midrule
Pascal (Thm.~\ref{thm:pascal}, \ref{thm:pascal-col}) & six points on $x^2+y^2=z^2$, 15 vertex pairs distinct; Pascal points pairwise distinct\\
Pappus (Cor.~\ref{cor:pappus}) & $A,C,E$ on $y=0$, $B,D,F$ on $x=0$; Pappus points pairwise distinct\\
Brianchon (Thm.~\ref{thm:brianchon}) & six tangents of the unit circle; diagonals pairwise distinct; common point $(7:-9:1)$\\
Prop.~\ref{prop:identity}--Rem.~\ref{rem:standard} & $\pi=-51$ at an explicit point; $c\ne0$, $Q\ne0$, $M\ne0$ symmetric; vertices and Pascal points distinct\\
Lemmas~\ref{lem:collinear}, \ref{lem:incidence}, \ref{lem:onconic} & a collinear and a non-collinear triple; distinct points and lines; a point on and a point off the conic\\
\bottomrule
\end{tabular}
\end{center}
\lean{l14\_apply\_eq\_coeffs} is an identity with no hypotheses and needs no witness. No
Mathlib-only statement turned out false or needed an extra hypothesis. Tangency of the six Brianchon
lines in the witness was checked numerically, not in Lean. The example in Section~\ref{sec:correct}
showing that $AB\ne DE$ cannot be dropped is not formalised. With the Layer-14 file added, the gate
checks 59 theorems (32 + 27) and passes, and the tampering suite rejects all six defects.
With the Layer-15 file added as well, the gate checks 147 theorems (59 + 88) and passes, and the
tampering suite again rejects all six defects. In that commit the GNU-only \texttt{mktemp} call in
\texttt{gate.sh} is replaced, so the gate also passes with the BSD tools of stock macOS;
\texttt{tests/tamper.sh} still needs GNU \texttt{sed}.

\paragraph{Layer 15: comparison with the classical statements.} Another AI agent collected the
classical statements from the literature, with citations (Pascal's \emph{Essay pour les coniques},
Brianchon's 1806 memoir and the modern projective form in~\cite{Coxeter}, Ch.~9). These were Pascal's theorem, its
five-point (tangent) form, Pappus's theorem, the converse (Braikenridge--Maclaurin), Brianchon's
theorem for a nondegenerate conic in characteristic $\ne2$, its polarity duality with Pascal, and its
pentagon form. A writer AI model was given only these statements; it saw neither the code nor this
paper, and its transcript was audited (a transcript audit in the verification repository). It wrote them in Lean with Mathlib's \lean{Projectivization},
with lines as $2$-dimensional subspaces of $K^3$, and with its own definitions of a nondegenerate
conic (``not equivalent to a form in fewer variables'', with at least one point), of the tangent at
$A$ (a line through $A$ that meets the conic in no other point) and of the polar line. This gave 20
statements, 12 theorems and 8 well-formedness obligations, all type-checked with \lean{sorry}. A
comparer then proved each one, copied verbatim, from the library, with no added hypothesis and no
change to the library, and classified it. All 20 are proved;
\lean{\#print axioms} reports only the three standard axioms for each:
\begin{itemize}
\item \emph{Covered} (6): Pascal's theorem (\lean{cov\_Pascal1\_a}, from \lean{pascal'})
and Pappus's theorem (\lean{cov\_Pascal2\_b}, from \lean{pappus}; the library file inlines the
one-line proof of \lean{pappus}, \lean{pascal'} applied to \lean{linePair}, so that the tampering
test that replaces \lean{pappus} is still rejected by the axiom check), with their three
obligations and the obligation that $XY$ is a line. In both cases the library statement is strictly more general. It allows any nonzero form
(resp.\ $a=b$ and vertices at $a\cap b$) and requires only distinct consecutive vertices and distinct
opposite sides.
\item \emph{Narrower} (1): Brianchon's theorem for six distinct tangent lines. Missing from the
library: the link between tangent lines and the dual conic, which Theorem~\ref{thm:brianchon} takes as
a hypothesis. NEW PROOF \lean{cov\_Brianchon1\_a} (characteristic $\ne2$) closes
the gap through the polarity of the conic, applied to \lean{pascal'}. It does not use
\lean{brianchon}.
\item \emph{Out of scope of the paper, now proved} (13), all NEW PROOFS:
  \begin{itemize}
  \item a nondegenerate conic contains no three collinear points, in any characteristic
  (\lean{cov\_Setting\_noThreeCollinear});
  \item the five-point form with the tangent at $A$, in any characteristic
  (\lean{cov\_Pascal2\_a}, with its obligation). The proof applies Lemma~\ref{lem:poly}
  over $K[X]$ and specialises $X=0$;
  \item the converse: six distinct points, no three collinear, whose Pascal points are collinear, lie
  on a nondegenerate conic, in any characteristic (\lean{cov\_Pascal3\_a}, with its
  obligation). The proof uses Proposition~\ref{prop:identity};
  \item the polarity (\lean{cov\_Brianchon2\_a\_}$X$ for $X$ = \lean{polarIsLine}, \lean{tangent},
  \lean{vertex}, \lean{incidence}, \lean{diagonal}). The tangent equals the polar in any characteristic; the others are stated
  and proved in characteristic $\ne2$;
  \item the obligation of Brianchon's theorem (opposite vertices are distinct), and the pentagon form
  with its obligation (\lean{cov\_Brianchon2\_b}).
  \end{itemize}
\item \emph{Writer statement false as written}: none.
\end{itemize}
The Layer-15 file reproduces the writer's definitions and statements verbatim. Relative to the
scratch file, the library version changes the namespace, header and imports, removes two linter
options, and inlines the proof of \lean{pappus} in \lean{cov\_Pascal2\_b}, as described above.


\section{Limitations and further work}

\begin{itemize}
\item \emph{Tangents.} The projective statement requires consecutive vertices to be distinct.
The limiting cases, in which a side is replaced by the tangent to the conic, are not part of the
main development, since Mathlib has no definition of the tangent line to a conic. With the
Layer-15 writer's definition (a line through $A$ meeting the conic only at~$A$), the five-point form
and the pentagon form of Brianchon's theorem are proved in the Layer-15 file
(Section~\ref{sec:l1415}), for nondegenerate conics only. Other limiting cases are not formalised.
\item \emph{The affine and Euclidean statements.} Wiedijk's list uses an affine or circle statement.
Deriving one from Theorem~\ref{thm:pascal-col} needs a bridge between Mathlib's affine notion
\lean{Collinear} in $\mathbb{R}^2$ and \lean{IsCollinear} in $\PP^2(\mathbb{R})$. We have not
formalised it.
\item \emph{Converse.} The converse of Pascal's theorem, that a hexagon whose opposite sides meet in
collinear points lies on a conic, follows from the same identity. It is not part of the main development. The Layer-15 file proves that six distinct points, no three collinear, whose Pascal points are
collinear lie on a nondegenerate conic in the writer's sense
(\lean{cov\_Pascal3\_a} in namespace \lean{PascalHexagon.Classical}, Section~\ref{sec:l1415}).
\item \emph{Tangency and the dual conic.} Formalising the correspondence between tangent lines
and the dual conic would make Theorem~\ref{thm:brianchon} a statement about tangent lines directly.
The tangent-line form itself is proved in the Layer-15 file through the polarity (characteristic
$\ne2$), but the correspondence with the dual conic is not.
\item \emph{Mathlib.} Contributing the result to Mathlib would require a general notion of conic
and faster proofs of Proposition~\ref{prop:identity}.
\end{itemize}

\section*{Acknowledgements and use of AI}
The formal proofs, the checking scripts and a draft of this paper were prepared with the help of
large-language-model AI assistants. A different model reviewed the formalisation. The author checked each formal statement in
Table~\ref{tab:thms} against its informal counterpart by hand. The Layer-14 and Layer-15 files were written by AI
agents and are not part of the author's hand check. The Layer-15 statements were written by a model that saw neither the code nor this paper. Every proof is
checked by the Lean kernel, and the author takes responsibility for all content.

\begin{thebibliography}{99}
\bibitem{AtlasLean} Meta Platforms, Inc., \emph{atlas-lean}, file
\path{v1/Atlas/AlgebraicGeometryI/code/PascalTheorem.lean},
\url{https://github.com/facebookresearch/atlas-lean}, accessed 28 September 2026.
\bibitem{Coghetto} R.~Coghetto, Pascal's theorem in real projective plane, \emph{Formalized
Mathematics} 25 (2017), no.~2, 107--119.
\bibitem{Coxeter} H.~S.~M.~Coxeter, \emph{Projective Geometry}, 2nd ed., Springer, 1987.
\bibitem{HOLLight} J.~Harrison, \texttt{100/pascal.ml}, HOL~Light repository,
\url{https://github.com/jrh13/hol-light}, accessed 28 September 2026.
\bibitem{LeanGenius} R.~J.~Walters, \emph{lean-genius}, file \texttt{proofs/Proofs/PascalsHexagon.lean},
\url{https://github.com/rjwalters/lean-genius}, accessed 28 September 2026.
\bibitem{MagaudNarboux} N.~Magaud, J.~Narboux and P.~Schreck, Formalizing projective plane geometry
in Coq, in \emph{Automated Deduction in Geometry (ADG 2008)}, Lecture Notes in Artificial
Intelligence 6301, Springer, 2011, 141--162.
\bibitem{Mathlib100} The mathlib Community, \texttt{docs/100.yaml}, Mathlib repository,
\url{https://github.com/leanprover-community/mathlib4}.
\bibitem{PlbyLean} \emph{lean-proofs}, file \texttt{src/latest/HundredTheorems/Theorem28.lean},
\url{https://github.com/plby/lean-proofs}, accessed 28 September 2026.
\bibitem{RichterGebert} J.~Richter-Gebert, \emph{Perspectives on Projective Geometry},
Springer, 2011.
\bibitem{Wiedijk} F.~Wiedijk, Formalizing 100 theorems, \url{https://www.cs.ru.nl/~freek/100/},
accessed 28 September 2026.
\end{thebibliography}

\end{document}
